%% sections/abstract.tex

\begin{abstract}
3D Gaussian Splatting (3DGS) has emerged as a transformative approach for
real-time radiance field rendering, combining the representational power of
explicit 3D primitives with differentiable rasterization.
Since the seminal work of Kerbl et~al.\ (2023), the field has experienced
explosive growth, with over \todo{N} papers published within two years.
However, this rapid expansion has led to a fragmented landscape where
contributions are difficult to compare and contextualize.

In this survey, we present a comprehensive taxonomy of 3DGS research,
organized along three orthogonal axes: \textbf{(1)~Task Domain}
(reconstruction, dynamic/4D, editing, generation, compression),
\textbf{(2)~Module Focus} (representation, optimization, rendering,
initialization, densification), and \textbf{(3)~Technique}
(anchoring, compression, anti-aliasing, regularization, level-of-detail,
feature augmentation).
We further introduce the \textbf{Drum Tower Dataset}, a new real-world
benchmark featuring Chinese Dong minority architecture with intricate
wooden structures and multi-scale geometry.
Our extensive experiments evaluate \todo{12} representative methods across
\todo{5} datasets, providing detailed analysis of quality-speed trade-offs,
memory consumption, and geometric accuracy.
Finally, we identify open challenges and promising future directions,
including scene-level understanding, physics-aware representations,
and edge-device deployment.

The project page and benchmark code are available at
\url{https://github.com/autumn119/3DGS-Survey-Taxonomy-Benchmark}.
\end{abstract}

\begin{IEEEkeywords}
3D Gaussian Splatting, Novel View Synthesis, Neural Rendering,
Radiance Fields, Survey, Benchmark, Taxonomy
\end{IEEEkeywords}
